// 323111047

\section{Introducción}

Planteamiento del problema. El reto técnico de esta práctica consiste en desarrollar un sistema de atención mixto mediante una interfaz de línea de comandos en C. El problema se centra en la implementación simultánea de dos tipos de estructuras de datos lineales: una cola doble para gestionar usuarios con prioridad VIP y una cola circular para usuarios convencionales. El reto principal radica en programar la lógica de atención de modo que el sistema respete siempre la prioridad de la estructura VIP sobre la circular, además de gestionar correctamente los punteros y el reciclaje de posiciones en el caso de la cola circular.

\vspace{0.4cm}

Motivación. La importancia de resolver este problema se basa en la necesidad de diseñar sistemas de gestión de recursos que imiten escenarios de la vida real, como la atención en bancos o aeropuertos, donde la segmentación por niveles de servicio es vital. Estudiar las colas dobles y circulares permite al estudiante entender cómo optimizar el uso de la memoria, evitando el desperdicio de espacios vacíos en arreglos, y cómo manejar flujos de datos con jerarquías.

\vspace{0.4cm}

Objetivos. El propósito fundamental de esta práctica es lograr la implementación funcional de un menú que permita registrar, atender y visualizar usuarios en dos estructuras distintas de forma organizada. Se busca que el alumno demuestre dominio en la manipulación de los frentes y finales de una cola circular, así como en las operaciones de inserción y extracción de una cola doble. Finalmente, se pretende validar que el programa cumpla estrictamente con la regla de prioridad, atendiendo a todos los usuarios VIP antes de pasar a la cola normal, comprobando así la correcta integración de la lógica de control en una aplicación de software.